%
% sinxx.tex -- Graph von sin(x)/x
%
% (c) 2021 Prof Dr Andreas Müller, OST Ostschweizer Fachhochschule
%
\documentclass[tikz]{standalone}
\usepackage{amsmath}
\usepackage{times}
\usepackage{txfonts}
\usepackage{pgfplots}
\usepackage{csvsimple}
\definecolor{darkred}{rgb}{0.8,0,0}
\definecolor{darkgreen}{rgb}{0,0.6,0}
\usetikzlibrary{arrows,intersections,math}
\input{sinxxdata.tex}
\begin{document}
\def\skala{1}
\def\s{1.5}
\begin{tikzpicture}[>=latex,thick,scale=\skala,
declare function = {
	f(\X) = \s*sin(180*\X)/\X;
}]

\begin{scope}
	\clip \stammfunktion -- (10.1,-1.2) -- (0,-1.2) -- cycle;
	\foreach \x in {1,3,...,9}{
		\fill[color=blue!5] (\x,-2) rectangle ++(1,10);
		\fill[color=darkred!5] (\x,-2) rectangle ++(-1,10);
	}
	\fill[color=darkred!5] (10,-2) rectangle ++(1,10);
\end{scope}

\fill[color=darkred!20] plot[domain=-1:1,samples=41]
	({\x},{f(\x)}) -- cycle;
\foreach \X in {2,4,6,8}{
	\fill[color=darkred!20] plot[domain=\X:{\X+1},samples=40]
		({\x},{f(\x)}) -- cycle;
}
\foreach \X in {2,4,6,8,10}{
	\fill[color=blue!20] plot[domain={\X-1}:{\X},samples=40]
		({\x},{f(\x)}) -- cycle;
}

\node[color=white] at (0.4,1) {$+$};
\node[color=white] at (2.47,0.3) {$+$};
\node[color=white] at (4.48,0.17) {$+$};
\node[color=white] at (6.5,0.11) {$+$};
\node[color=white] at (8.5,0.09) {$+$};

\node[color=white] at (1.5,-0.4) {$-$};
\node[color=white] at (3.5,-0.2) {$-$};
\node[color=white] at (5.5,-0.13) {$-$};
\node[color=white] at (7.5,-0.10) {$-$};
\node[color=white] at (9.5,-0.08) {$-$};

\begin{scope}
	\clip (-1.05,-1.2) rectangle (10.1,5);
	\draw[color=darkred,line width=1.2pt]
		plot[domain=-1.1:10.2,samples=600] ({\x},{f(\x)});
\end{scope}

\begin{scope}
	\clip \stammfunktion -- (10,-1.2) -- (0,-1.2) -- cycle;
	\foreach \x in {1,...,10}{
		\draw[line width=0.2pt] (\x,5) -- (\x,-1.2);
	}
\end{scope}

\coordinate (S) at (1,4.5);

\node[color=darkred] at (S) [right] {$\displaystyle f(x)=\frac{\sin x}{x}$};
\draw[color=darkred,line width=0.2pt] (S) -- (0.25,{f(0.25)});

\draw[color=darkgreen,line width=1.2pt] \stammfunktion;
\node[color=darkgreen] at (5,3) {$\displaystyle F(x) = \int_0^x \frac{\sin t}{t}\,dt$};

\draw[->] (-1.1,0) -- (10.3,0) coordinate[label={$x$}];
\draw[->] (0,-1.2) -- (0,5.0) coordinate[label={right:$y$}];

\draw (-1,-0.05) -- ++(0,0.1);
\node at (0,-1.2) [below] {$0\mathstrut$};
\draw (1,-0.05) -- ++(0,0.1);
\node at (1,-1.2) [below] {$\pi\mathstrut$};
\foreach \x in {2,...,10}{
	\draw (\x,-0.05) -- ++(0,0.1);
	\node at (\x,-1.2) [below] {$\x\pi\mathstrut$};
}

\end{tikzpicture}
\end{document}

